% Source: https://artofproblemsolving.com/community/c3h2297889_proving_theorems
% Stable identity: aops-topic-2297889-q02
In the figure, $\angle1,\angle2$ and $\angle3,\angle4$ are two pairs of supplementary angles, drawn as linear pairs.

\begin{asy}
usepackage("amsmath");
size(6cm);
defaultpen(fontsize(10pt));
pair O=(0,0); draw((-1.6,0)--(1.6,0),linewidth(0.9)); draw(O--1.5*dir(55),linewidth(0.9));
draw(arc(O,0.4,0,55)); draw(arc(O,0.5,55,180));
label("$1$",0.62*dir(27),fontsize(9pt)); label("$2$",0.72*dir(118),fontsize(9pt));
pair O2=(4,0); draw((2.4,0)--(5.6,0),linewidth(0.9)); draw(O2--O2+1.5*dir(125),linewidth(0.9));
draw(arc(O2,0.4,125,180)); draw(arc(O2,0.5,0,125));
label("$3$",O2+0.62*dir(152),fontsize(9pt)); label("$4$",O2+0.72*dir(62),fontsize(9pt));
\end{asy}

\begin{tabular}{ll}
\textbf{Statements}&\textbf{Reasons}\\\hline
$\angle1,\angle2$ supplementary; $\angle3,\angle4$ supplementary; $\angle2\cong\angle4$&Given\\
$m\angle1+m\angle2=180^\circ$, $m\angle3+m\angle4=180^\circ$&Definition of supplementary angles\\
$m\angle1+m\angle2=m\angle3+m\angle4$&Transitive property of equality\\
$m\angle2=m\angle4$&Definition of congruent angles\\
$m\angle1=m\angle3$&Subtraction property of equality\\
$\angle1\cong\angle3$&Definition of congruent angles
\end{tabular}

For two angles supplementary to the same angle, take $\angle4=\angle2$. $\blacksquare$
